% ============================================================
%  Reeds Scan  —  Methodology
% ============================================================

\subsection{Reeds Scan}

Two radar scans of the reed bed were acquired to test detection performance
under different levels of foliage occlusion.
The first scan was conducted on 11~March~2026, with the corner reflectors
elevated on stands so that they were only partially occluded by the upper
reed stems.
The second scan was conducted on 8~April~2026, with the corner reflectors
placed directly at ground level behind the reeds, resulting in heavy occlusion
with the reflectors barely visible to the naked eye.
Keeping the same radar scan parameters across both sessions allowed a direct
comparison of how occlusion severity affects detection performance.

\subsubsection*{Scan Procedure}

The scan parameters used on both days are given in Table~\ref{tab:scan_params}.
Both scans were conducted under sunny, slightly windy conditions at
approximately 27$^\circ$C with no cloud cover.
On both days, two corner reflectors were placed side by side behind the reeds,
with the radar positioned at the same location for consistency.

In the March scan, the reflectors were elevated on stands above the densest
part of the reed stems, reducing the effective path length through vegetation.
In the April scan, the reflectors were placed on gravel at ground level,
maximising the path length through the full reed height.
Figure~\ref{fig:reeds_setup_march} and Figure~\ref{fig:reeds_setup_april}
show the physical setup for each session.

\subsubsection*{Behind-Foliage Region Definition}

The detection analysis was spatially constrained to the region between the
front face of the reeds and the corner reflector.
This definition was chosen because it includes the vegetation returns as part
of the clutter, which gives a more realistic false positive rate than using
only the narrow gap immediately behind the foliage.
Using the back face of the reeds as the boundary was considered but rejected,
as the resulting region contained only 8--10 points, making any false positive
rate statistically meaningless.
The reed front face was identified at a range of 4.1~m in March and 3.0~m in
April from the radar.
All returns in this window were included in the clutter floor and false
positive calculations.

\subsubsection*{Detection Threshold Approaches}

Two complementary threshold approaches were used to analyse detection
performance, each serving a different analytical purpose.

The first approach used a fixed threshold derived from the control scan
clutter floor ($-40.1$~dB).
Applying the same threshold to both foliage scans enables direct comparison
between March and April on a common reference, and allows vegetation loss to
be computed relative to free-space conditions.

The second approach estimated the clutter floor independently for each scan
from the returns within the behind-foliage region.
This is more representative of how a threshold would be set in a real
deployment, where no prior free-space reference is available.
The trade-off is that it makes direct comparison between scans harder, since
each scan operates against a different baseline.

Both approaches detected the CR successfully in all scans.
The fixed threshold approach is used for the vegetation loss calculations,
and the scan-specific approach is used for the ROC curve comparison.

\subsubsection*{Vegetation Loss Calculation}

Vegetation loss was estimated by comparing the measured CR peak amplitude
in each foliage scan against the expected free-space amplitude at the same
range.
The expected free-space amplitude was derived by applying the free-space
propagation model established from the control scan to the range of each CR.
The vegetation loss is then:
\[
  L_{\text{veg}} = A_{\text{free-space}} - A_{\text{measured}}
\]
where $A_{\text{free-space}}$ is the predicted amplitude in the absence of
vegetation and $A_{\text{measured}}$ is the actual peak amplitude observed
in the foliage scan.


% ============================================================
%  Reeds Scan  —  Results
% ============================================================

\subsubsection*{Point Cloud Visualisation}

Figure~\ref{fig:reeds_pointcloud} shows the 3D radar point clouds for the
March and April scans coloured by amplitude.
In the March scan, both CRs are visible as high-amplitude clusters standing
above the surrounding reed returns.
In the April scan, the return pattern changes considerably — April CR1
produces an unusually strong return while CR2 is largely suppressed.
The dominant factor driving this difference is the height of the reflectors
relative to the reed canopy, which determines the effective path length through
vegetation and is examined quantitatively below.

\subsubsection*{Detection Results}

\paragraph{Scan-specific thresholds.}
For the March scan, the behind-foliage clutter floor was $-31.8$~dB
($\sigma = 22.0$~dB).
The CR cluster at a range of 10.92~m peaked at 18.0~dB, giving an SNR of
49.8~dB above the local clutter floor, with 118 CR points detected.

For the April scan, the clutter floor was $-33.0$~dB.
April CR1, at a range of 13.88~m, peaked at 53.0~dB giving an SNR of
86.0~dB, with 261 CR points detected.
April CR2 peaked at only $-11.0$~dB, falling below the clutter floor and
remaining undetectable — consistent with total occlusion by the dense reed
stems at ground level.

The ROC curve comparison in Figure~\ref{fig:reeds_roc_comparison} shows that
March (partial occlusion) achieves higher true positive rates across most
threshold values.
Both scans successfully detected the CR at their respective clutter floors,
confirming that even partial occlusion does not prevent detection when the
reflector is elevated above the densest vegetation layer.

\paragraph{Fixed control threshold.}
At the fixed threshold of $-40.1$~dB, the March scan achieved a true positive
rate (TPR) of 0.93 and a false positive rate (FPR) of 0.67.
The April scan achieved a TPR of 0.57 with the same FPR of 0.67.
The identical FPR across both scans indicates that the behind-foliage clutter
distribution was similar in both sessions, meaning the vegetation structure
did not change meaningfully between March and April.
The lower TPR for April is consistent with heavier occlusion reducing the
number of detectable CR returns above the threshold.
The full detection summary is given in Table~\ref{tab:control_referenced_summary}.

\subsubsection*{Vegetation Attenuation Estimate}

For the March scan, the expected free-space amplitude at 10.92~m was 76.3~dB.
The measured CR peak was 18.0~dB, giving a vegetation loss of:
\[
  L_{\text{March}} = 76.3 - 18.0 = 58.3~\text{dB}
\]

For the April scan, the expected free-space amplitude at 13.88~m was 72.2~dB.
The measured CR1 peak was 53.0~dB, giving a vegetation loss of:
\[
  L_{\text{April}} = 72.2 - 53.0 = 19.2~\text{dB}
\]

Figure~\ref{fig:cr_amplitude_range} plots the CR peak amplitudes for both
scans against range, with the detection threshold of $-29.2$~dB shown as a
dashed reference line.
All detected CRs sit well above this threshold.

The April vegetation loss (19.2~dB) is considerably lower than March (58.3~dB),
despite the April reflectors being more heavily occluded.
This is counterintuitive and is addressed below.

\paragraph{April anomaly — multipath contribution.}
The unexpectedly high return from April CR1 is attributed to multipath
reflections from the gravel surface beneath the reflector.
In April, the CRs were placed on gravel, which is harder and more reflective
than the grass used in March.
At ground level, the radar signal can bounce off the gravel and reach the CR
via an indirect path, constructively combining with the direct return and
boosting the measured amplitude.
In March, the CRs were elevated on grass, which absorbs more energy and
produces less ground-bounce multipath.
This surface material difference likely accounts for a significant portion
of the 35~dB amplitude difference between March and April CR1, independent
of vegetation loss.
The April vegetation loss of 19.2~dB should therefore be treated as a lower
bound on the true path loss through the reeds at ground level.

\subsubsection*{Amplitude Distributions}

Figure~\ref{fig:reeds_amplitude_hist} shows the amplitude distributions for
both scans.
The March distribution is broad and relatively flat, centred around $-25$~dB,
with the CR peaks at approximately 18~dB appearing as a clear spike at the
far right — well separated from the vegetation clutter.

The April distribution is more compact and centred around $-28$~dB.
April CR2 ($-11.0$~dB) sits within the upper tail of the clutter distribution,
explaining why it is undetectable.
April CR1 (53.0~dB) is an extreme outlier far beyond the clutter body,
consistent with the multipath-boosted return discussed above.

The broader March distribution reflects greater spatial variability in returns
from partially-occluded vegetation, while the more compact April distribution
suggests a more uniform clutter background at ground level.

\subsubsection*{LiDAR Scans}

The LiDAR scan of the reed bed confirmed that the 905~nm laser cannot
penetrate the vegetation.
Figure~\ref{fig:reeds_lidar} shows a high density of surface returns from
the front face of the reed bed, with the point count dropping sharply beyond
this surface.
Any corner reflectors placed behind the reeds produce no detectable LiDAR
return.
This is in direct contrast to the radar, which successfully detected the CR
through the reeds in all tested configurations.
The result demonstrates the fundamental advantage of millimetre-wave radar
over optical LiDAR for concealed-target detection in dense vegetation.
